This thesis presents measurements of the azimuth-independent amplitude $A_0$ and the cosine moments $A_{UU}^{\cos\phi_h}$ and $A_{UU}^{\cos 2\phi_h}$ for semi-inclusive $\pi^{+}$ production from unpolarized hydrogen at a beam energy of $10.604\,\mathrm{GeV}$. The measurements cover $452$ retained $Q^{2}$--$y$--$z$--$P_T$ bins. The nominal results use the spline-weighted analysis after acceptance, $RC$, and $BC$ corrections (\Cref{Cosine Moment Extractions,sys-mod-spline}). Over the $452$ bins the median of $A_{UU}^{\cos\phi_h}$ is $-0.082$ and the median of $A_{UU}^{\cos 2\phi_h}$ is $-0.0050$ (\Cref{tab:CahnGaussian}). $A_{UU}^{\cos\phi_h}$ is consistent with zero near the lowest $P_{T}$ and becomes negative as $P_{T}$ increases. $A_{UU}^{\cos 2\phi_h}$ does not follow that slope and remains near zero. The three moments are shown against $z$ and against $P_{T}$ in \Crefrange{fig:Moments_Spline_z_A_sys}{fig:Moments_Spline_pT_C_sys}.

The comparisons with published $\pi^{+}$ moments on hydrogen are given in \Cref{Comparisons to World Data,tab:WorldOverlap,fig:WorldMomentsVsPT}. The median absolute difference from the Harrison e1-f moments is $0.037$ in $A_{UU}^{\cos\phi_h}$ and $0.037$ in $A_{UU}^{\cos 2\phi_h}$, over $29$ points in $17$ bins.~\cite{HarrisonThesis} The median absolute difference from the Osipenko moments at $5.75\,\mathrm{GeV}$, after the depolarization conversion in that section, is $0.105$ and $0.024$, over $24$ points in $18$ bins.~\cite{Osipenko:2008rv} The $\cos 2\phi_h$ entry uses $23$ points, because one of those bins does not determine the moment. Differences involving the same present bin share its measurement uncertainty and are therefore correlated. The fall of $A_{UU}^{\cos\phi_h}$ with $P_{T}$ has the same sign in both comparisons. At the lower beam energy those measurements sit at higher $y$, so the depolarization factors in \Cref{eq:Cosphi_Moment,eq:Cos2phi_Moment} are smaller. Dividing each moment by its own factor reduces the median absolute $\cos\phi_h$ difference to $0.011$ for Harrison and $0.029$ for Osipenko, and it raises the $\cos 2\phi_h$ differences to $0.072$ and $0.048$. The change of sign with $P_{T}$ in the Harrison $\cos 2\phi_h$ points is not removed by that factor. Hall~C and HERMES have also published cosine modulations relevant to this measurement.~\cite{HallCSIDIS:2025plr,Airapetian:2012yg} A quantitative comparison with the released Hall~C data has not yet been completed. Extraction and validation of the corresponding cosine coefficients are planned as this work progresses toward publication. The quantitative comparison with HERMES remains deferred because tabulated numerical values of its fitted coefficients were not obtained for the present comparison. The cited COMPASS measurement provides experimental context, with a different target, hadron selection, and $x_B$ coverage.~\cite{Adolph:2014pwc}

The theory comparisons evaluate published formulas at the $452$ bin centers (see \Cref{Comparisons to Theory Models}). The published-width Cahn term (see ~\cite{Anselmino:2005nn,Barone:2015cgh}), with $\langle k_{\perp}^{2}\rangle = 0.25\,\mathrm{GeV}^{2}$ and $\langle p_{\perp}^{2}\rangle = 0.20\,\mathrm{GeV}^{2}$, is more negative overall than the measured central values. Its median is $-0.217$, and the root-mean-square difference from the measurement is $0.183$. The kinematic $\cos 2\phi_h$ term from the same Gaussian integral has median $0.0067$, against the measured median $-0.0050$, and the root-mean-square difference is $0.033$. The Boer--Mulders--Collins term taken from Bastami et al.~\cite{Bastami:2018xqd} is negative in every bin. Its median is $-0.0065$, and the root-mean-square difference from the measurement is $0.037$ (\Cref{tab:BoerMuldersBastami}). The Cahn and Boer--Mulders--Collins contributions are evaluated separately as model benchmarks. These comparisons do not separate the measured moments into individual mechanisms or determine the Boer--Mulders function or quark orbital angular momentum.

The $P_T^{2}$ dependence of $A_{0}$ is used to extract the effective Gaussian width $\kappa_T^{2}(z)$ over the fitted range (\Cref{Transverse Momentum Extraction,tab:MeanKT2,fig:LnA0_all,fig:kT2_all}). The extracted widths are presented as functions of $z$ in each $Q^{2}$--$y$ bin. Their central values range from $0.160\,\mathrm{GeV}^{2}$ to $0.304\,\mathrm{GeV}^{2}$. The point uncertainties and the fit weights use only the systematic uncertainties on $A_{0}$, and the uncertainties on the widths are the fit uncertainties. One fit, in $Q^{2}$--$y$ bin 15 at $z=0.555$, has a fit uncertainty at least half the width, as marked in \Cref{tab:MeanKT2}. That fit is not unsuccessful.

The systematic uncertainties were evaluated from variations of the analysis choices described in \Cref{Systematic Uncertainties}. Their contributions are calculated for each measurement and combined to obtain its systematic uncertainty. \Cref{tab:SystematicUncertainties} summarizes the effects of the individual sources across the retained kinematic bins; the uncertainties shown with the measurements are determined for each bin.

These multidimensional measurements provide inputs for phenomenological studies of transverse-momentum dynamics in unpolarized SIDIS. The present comparisons do not separate the physical mechanisms contributing to the cosine moments. Such a separation, or a determination of separate quark and fragmentation transverse-momentum widths, would require additional model assumptions and/or complementary measurements. The available $Q^2$ range is not used here to study scale dependence.
